\begin{tcolorbox}[colback=red!3!white,colframe=red!50!black,boxrule=0.8pt,title={Correction notes},fonttitle=\bfseries]
\textbf{Corrections from the calculation check (30~Sept.~2026).} The following places are corrected or annotated in the text; each carries a dated note with the previous value:
\begin{itemize}
\item $\alpha$ error $0.006\,\% \to 0.0005\,\%$ (three places); reconstruction from mass ratios: $0.02747 \to 0.02860$, result $0.00730 \to 0.00759$ ($1/\alpha = 131.7$).
\item Koide value of the T0 masses: $Q_{\text{T0}} = 0.666664 \pm 0.000005$, $\Delta Q = 0.00003\,\%$ $\to$ $Q_{\text{T0}} = 0.667742$, $\Delta Q = +0.161\,\%$ (also comparison table and overview); main theorem and \enquote{Purest $\xi$-Form} annotated. Second line of the $Q$ calculation: superfluous $\cdot v$ in the denominator removed.
\item Empirical accuracy: $\Delta Q < 0.00003\,\% \to |\Delta Q| \approx 0.0003\,\%$ ($\pm 0.0008\,\%$ from $m_\tau$); Addendum P1 annotated.
\item $G$ precision (note of 1~Oct.~2026): $C_{\text{conv}}$ is the SI conversion of the single-anchor chain and additionally contains its small residual (Doc.~383); the agreement to $0.003\,\%$ does not count as an independent precision result (Doc.~190, R141).
\end{itemize}
\end{tcolorbox}

\section*{Abstract}
		We prove that the Koide formula for lepton masses is not an independent empirical relation, but a mathematical consequence of the geometric constant $\xi = \frac{4}{3} \times 10^{-4}$ from the T0 theory. The quantum ratios $(r,p)$ of the T0-Yukawa formula $m = r \cdot \xi^p \cdot v$ automatically generate the Koide symmetry $Q = \frac{2}{3}$ without additional parameters or fractal corrections.

	
	\section{The Koide Formula}
	
	The relation discovered by Yoshio Koide in 1981 connects the masses of the charged leptons:
	
	\begin{equation}
		Q = \frac{m_e + m_\mu + m_\tau}{\left( \sqrt{m_e} + \sqrt{m_\mu} + \sqrt{m_\tau} \right)^2} = \frac{2}{3}
		\label{eq:koide}
	\end{equation}
	
	With the PDG 2024 masses ($m_\tau = 1776.93 \pm 0.09$~MeV) this formula achieves an experimental accuracy of $|\Delta Q| \approx 0.0003\%$; the uncertainty of $m_\tau$ alone gives $\pm 0.0008\%$.
	
	\section{T0-Yukawa Formula}
	
	In the T0 theory, particle masses arise from:
	
	\begin{equation}
		m = r \cdot \xi^p \cdot v
		\label{eq:t0yukawa}
	\end{equation}
	
	with Higgs VEV $v = 246$ GeV and $\xi = \frac{4}{3} \times 10^{-4}$.
	
	\subsection{Lepton Parameters}
	
	\begin{table}[h]
		\centering
		\adjustbox{max width=\textwidth}{%
\begin{tabular}{lccc}
			\toprule
			\textbf{Lepton} & \textbf{$r$} & \textbf{$p$} & \textbf{$m$ [GeV]} \\
			\midrule
			Electron & $\frac{4}{3}$ & $\frac{3}{2}$ & 0.000511 \\
			Muon & $\frac{16}{5}$ & $1$ & 0.1057 \\
			Tau & $\frac{25}{9}$ & $\frac{2}{3}$ & 1.7769 \\
			\bottomrule
		\end{tabular}%
}
		\caption{T0 Quantum Ratios of the Charged Leptons}
	\end{table}
	
	\section{Main Theorem}
	
	\begin{theorem}
		The Koide relation $Q = \frac{2}{3}$ is a direct mathematical consequence of the T0 exponents $(p_e, p_\mu, p_\tau) = \left(\frac{3}{2}, 1, \frac{2}{3}\right)$ and the associated ratios $(r_e, r_\mu, r_\tau) = \left(\frac{4}{3}, \frac{16}{5}, \frac{25}{9}\right)$.
	\end{theorem}
	
	\section{Proof via Mass Ratios}
	
	\subsection{Electron to Muon}
	
	\begin{beweis}
		\allowdisplaybreaks
		\begin{align}
			\frac{m_e}{m_\mu} &= \frac{r_e \cdot \xi^{p_e}}{r_\mu \cdot \xi^{p_\mu}} = \frac{\frac{4}{3} \cdot \xi^{3/2}}{\frac{16}{5} \cdot \xi^1} \\
			&= \frac{4}{3} \cdot \frac{5}{16} \cdot \xi^{1/2} = \frac{5}{12} \cdot \xi^{1/2} \\
			&= \frac{5}{12} \cdot \sqrt{1.333 \times 10^{-4}} \\
			&= \frac{5}{12} \cdot 0.01155 = 0.004813 \\
			&\approx \frac{1}{206.768} \quad \checkmark
		\end{align}
		
		\textbf{Experimental:} $\frac{m_e}{m_\mu} = 0.004836$ (PDG 2024)\\
		\textbf{Deviation:} $< 0.5\%$
	\end{beweis}
	
	\subsection{Muon to Tau}
	
	\begin{beweis}
		\begin{align}
			\frac{m_\mu}{m_\tau} &= \frac{r_\mu \cdot \xi^{p_\mu}}{r_\tau \cdot \xi^{p_\tau}} = \frac{\frac{16}{5} \cdot \xi^1}{\frac{25}{9} \cdot \xi^{2/3}} \\
			&= \frac{16}{5} \cdot \frac{9}{25} \cdot \xi^{1/3} = \frac{144}{125} \cdot \xi^{1/3} \\
			&= 1.152 \cdot (1.333 \times 10^{-4})^{1/3} \\
			&= 1.152 \cdot 0.05105 = 0.05881 \\
			&\approx \frac{1}{17.0} \quad \checkmark
		\end{align}
		
		\textbf{Experimental:} $\frac{m_\mu}{m_\tau} = 0.05947$ (PDG 2024)\\
		\textbf{Deviation:} $< 1.2\%$
	\end{beweis}
	
	\subsection{Electron to Tau}
	
	\begin{beweis}
		\begin{align}
			\frac{m_e}{m_\tau} &= \frac{r_e \cdot \xi^{p_e}}{r_\tau \cdot \xi^{p_\tau}} = \frac{\frac{4}{3} \cdot \xi^{3/2}}{\frac{25}{9} \cdot \xi^{2/3}} \\
			&= \frac{4}{3} \cdot \frac{9}{25} \cdot \xi^{5/6} = \frac{12}{25} \cdot \xi^{5/6} \\
			&= 0.48 \cdot (1.333 \times 10^{-4})^{5/6} \\
			&= 0.48 \cdot 0.000590 = 0.000283 \\
			&\approx \frac{1}{3534} \quad \checkmark
		\end{align}
		
		\textbf{Experimental:} $\frac{m_e}{m_\tau} = 0.0002876$ (PDG 2024)\\
		\textbf{Deviation:} $< 1.6\%$
	\end{beweis}
	
	\section{Direct Derivation of the Koide Relation}
	
	\subsection{Geometric Structure of the Exponents}
	
	The T0 exponents exhibit a fundamental symmetry:
	
	\begin{equation}
		p_e - p_\mu = \frac{3}{2} - 1 = \frac{1}{2}
	\end{equation}
	\begin{equation}
		p_\mu - p_\tau = 1 - \frac{2}{3} = \frac{1}{3}
	\end{equation}
	
	These generate the characteristic $\sqrt{m}$-dependencies of the Koide formula.
	
	\subsection{Calculation of $Q$}
	
	Substituting the T0 masses into equation \eqref{eq:koide}:
	
	\begin{align}
		Q &= \frac{r_e \xi^{p_e} v + r_\mu \xi^{p_\mu} v + r_\tau \xi^{p_\tau} v}{\left(\sqrt{r_e \xi^{p_e} v} + \sqrt{r_\mu \xi^{p_\mu} v} + \sqrt{r_\tau \xi^{p_\tau} v}\right)^2} \\
		&= \frac{r_e \xi^{3/2} + r_\mu \xi + r_\tau \xi^{2/3}}{\left(\sqrt{r_e} \xi^{3/4} + \sqrt{r_\mu} \xi^{1/2} + \sqrt{r_\tau} \xi^{1/3}\right)^2}
	\end{align}
	\textit{(Corrected on 30~Sept.~2026: previously the denominator of the second line still contained $\cdot v$ -- $v$ cancels completely in numerator and denominator, $Q$ is dimensionless.)}
	
	With the numerical values:
	\begin{align}
		Q_{\text{T0}} &= 0.667742 \\
		Q_{\text{Koide}} &= \frac{2}{3} = 0.666667 \\
		\Delta Q &= +0.161\%
	\end{align}
	\textit{(Corrected on 30~Sept.~2026: previously $Q_{\text{T0}} = 0.666664 \pm 0.000005$, $\Delta Q = 0.00003\%$ -- with $m_i = r_i \xi^{p_i} v$, $r = (4/3, 16/5, 25/9)$, $p = (3/2, 1, 2/3)$ one obtains $m = (0.505, 104.96, 1783.4)$~MeV and $Q_{\text{T0}} = 0.667742$, i.e.\ $0.161\%$ above $2/3$. With these $(r,p)$ the main theorem $Q = 2/3$ is not fulfilled numerically; the conclusions of this section hold only approximately ($Q \approx 2/3$).)}
	
	\section{Key Insight}
	
	\begin{folgerung}
		\textbf{The Koide formula is not an independent symmetry, but a direct manifestation of $\xi$.}
		
		\begin{itemize}
			\item The exponents $(3/2, 1, 2/3)$ generate the $\sqrt{m}$-structure
			\item The ratios $(4/3, 16/5, 25/9)$ are consistent with $Q \approx 2/3$
			\item No fractal corrections necessary
			\item No additional free parameters
			\item The geometric constant $\xi$ was implicitly already contained in the Koide formula
		\end{itemize}
	\end{folgerung}
	
	\section{Comparison: Empirical vs. T0 Derivation}
	
	\begin{table}[h]
		\centering
		\adjustbox{max width=\textwidth}{%
\begin{tabular}{lcc}
			\toprule
			\textbf{Aspect} & \textbf{Koide (1981)} & \textbf{T0 Theory} \\
			\midrule
			Free Parameters & 0 (empirical) & 1 ($\xi$) \\
			Basis & Observation & Geometry \\
			Accuracy & $\approx 0.0003\%$ & $0.16\%$ \\
			Explanation & None & $\xi$-Geometry \\
			Predictive Power & Only Leptons & All Particles \\
			\bottomrule
		\end{tabular}%
}
		\caption{Comparison of Approaches}
	\end{table}
	\textit{(Corrected on 30~Sept.~2026: previously $< 0.00003\%$ twice in the row \enquote{Accuracy} -- empirically $|\Delta Q| \approx 0.0003\%$ (PDG masses), T0 masses $Q_{\text{T0}} = 0.667742$, i.e.\ $0.16\%$.)}
	
	\section{Mathematical Significance}
	
	The T0 formula shows that:
	
	\begin{equation}
		Q = \frac{2}{3} \iff \text{Exponents form geometric series with base } \xi
	\end{equation}
	
	This explains:
	\begin{enumerate}
		\item Why $Q = 2/3$ and not another value
		\item Why the relation applies to exactly 3 generations
		\item Why square roots of masses (not masses themselves) are added
		\item The connection to Higgs-Yukawa coupling
	\end{enumerate}
	
	\section{Fine Structure Constant from Mass Ratios}
	
	\subsection{Direct T0 Derivation}
	
	The fine structure constant in the T0 theory:
	
	\begin{equation}
		\alpha = \xi \cdot \left(\frac{E_0}{1\,\text{MeV}}\right)^2 = \frac{4}{3} \times 10^{-4} \times (7.398)^2 = 0.007297
	\end{equation}
	
	where $E_0$ is derived from the lepton mass ratios, as shown in the following subsection.
	
	\textbf{Experimental:} $\alpha = \frac{1}{137.036} = 0.0072973525693$\\
	\textbf{Error:} $0.0005\%$ \textit{(Corrected on 30~Sept.~2026: previously $0.006\%$ -- $1/(\xi\cdot7.398^2) = 137.035$ against $137.036$ is $4.7\times10^{-6}$; likewise in the box and the overview below.)}
	
	\subsection{Reconstruction from Lepton Masses}
	
	\begin{beweis}
		The fine structure constant can be reconstructed from the mass ratios:
		
		\begin{equation}
			\alpha \propto \left(\frac{m_e}{m_\mu}\right)^{2/3} \times \left(\frac{m_\mu}{m_\tau}\right)^{1/2} \times \xi^{\text{const}}
		\end{equation}
		
		With the T0 ratios:
		\begin{align}
			\alpha_{\text{rekon}} &= \left(\frac{1}{206.768}\right)^{2/3} \times \left(\frac{1}{16.818}\right)^{1/2} \times 1.089 \\
			&= 0.02860 \times 0.2438 \times 1.089 \\
			&\approx 0.00759
		\end{align}
		\textit{(Corrected on 30~Sept.~2026: previously $0.02747$ and $\approx 0.00730$ -- $(1/206.768)^{2/3} = 0.02860$, the product is $0.00759$, i.e.\ $1/\alpha = 131.7$ ($+4\,\%$); the factor $1.089$ is not derived. The reconstruction therefore does not hit $\alpha$.)}
	\end{beweis}
	
	\textbf{Remarkable:} The exponents $(2/3, 1/2)$ are directly linked to the T0 exponent differences:
	\begin{itemize}
		\item $p_e - p_\mu = \frac{3}{2} - 1 = \frac{1}{2}$ appears in $\sqrt{m_\mu/m_\tau}$
		\item $p_\mu - p_\tau = 1 - \frac{2}{3} = \frac{1}{3}$ appears in $(m_e/m_\mu)^{2/3}$
	\end{itemize}
	
	\section{Hierarchy of $\xi$-Manifestations}
	
	The three fundamental constants arise from $\xi$ at different ``purity levels'':
	
	\subsection{Level 1: Mass Ratios (Koide Formula)}
	
	\begin{equation}
		Q = \frac{\sum m_i}{\left(\sum \sqrt{m_i}\right)^2} \quad \text{with} \quad m_i = r_i \xi^{p_i} v
	\end{equation}
	
	\begin{tcolorbox}[colback=green!5!white,colframe=green!75!black,title={Purest $\xi$-Form}]
		\textbf{Accuracy:} $\Delta Q < 0.00003\%$
		\textit{(Note of 30~Sept.~2026: With $m_i = r_i \xi^{p_i} v$ one obtains $Q = 0.667742$, i.e.\ $\Delta Q = +0.161\%$; empirically (PDG masses) $|\Delta Q| \approx 0.0003\%$. The reasoning \enquote{Why perfect} therefore does not hold numerically.)}
		
		\textbf{Why perfect:}
		\begin{itemize}
			\item Only ratios, no absolute scales
			\item $\xi$ appears only in exponent differences: $\xi^{p_i - p_j}$
			\item Higgs VEV $v$ cancels completely
			\item NO fractal corrections necessary
		\end{itemize}
	\end{tcolorbox}
	
	\subsection{Level 2: Fine Structure Constant}
	
	\begin{equation}
		\alpha = \xi \cdot E_0^2
	\end{equation}
	
	\begin{tcolorbox}[colback=blue!5!white,colframe=blue!75!black,title={Semi-pure $\xi$-Form}]
		\textbf{Accuracy:} $\Delta \alpha \approx 0.0005\%$
		
		\textbf{Why very good:}
		\begin{itemize}
			\item Requires an energy scale $E_0 = 7.398$ MeV, which is emergently derived from the mass ratios
			\item Direct $\xi$-coupling
			\item Small uncertainty due to $E_0$-calibration
		\end{itemize}
	\end{tcolorbox}
	
	\subsection{Level 3: Gravitational Constant}
	
	\begin{equation}
		G = \frac{\xi^2}{4m} = \frac{\xi^2}{4 \cdot \xi/2} = \frac{\xi}{2} = 6.667 \times 10^{-5} \quad \text{(in natural units)}
	\end{equation}
	
	With SI conversion: $G_{\text{SI}} = G_{\text{nat}} \times 1.00115 \times 10^{-6}\,\text{m}^3\text{kg}^{-1}\text{s}^{-2}$
	
	\textit{(Corrected on 29~Sept.~2026: previously $\xi^2/(4\cdot\xi/2) = \xi$, conversion factor $2.843\times10^{-5}$ and accuracy $0.5\%$ -- the fraction gives $\xi/2$; the factor matching $G_{\text{nat}} = \xi/2$ is $K_{\text{SI}} = 1.00115\times10^{-6}$ as in Doc.~028; the corpus formula $G = \xi^2/(4m_e)\cdot C_{\text{conv}}\cdot K_{\text{frak}}$ yields $6.6745\times10^{-11}$, i.e.\ $+0.003\%$ relative to CODATA.)} \textit{(Note of 1~Oct.~2026: $C_{\text{conv}}$ is the SI conversion of the single-anchor chain (Doc.~383); its numerical value $7.783\times10^{-3}$ additionally contains the small residual of this chain (chain value $7.58\times10^{-3}$). The form of $G$ and its order of magnitude follow from $\xi$ and the one anchor (numerically checked in Doc.~383, [K]); with the anchor $m_e$, $G$ lies at $-2.6\,\%$, with $v$ at $-0.46\,\%$ (Doc.~383). The agreement to $0.003\,\%$ arises only with this residual and does not count as an independent precision result (Doc.~190, R141).)}
	
	\begin{tcolorbox}[colback=yellow!5!white,colframe=orange!75!black,title={Complex $\xi$-Form}]
		\textbf{Accuracy:} $\Delta G \approx +0.003\%$
		
		\textbf{Why more difficult:}
		\begin{itemize}
			\item Requires Planck length $\ell_P = 1.616 \times 10^{-35}$ m, which is directly related to $\xi$ ($\ell_P \propto \sqrt{G} \propto \sqrt{\xi}$ in natural units)
			\item Complex SI units conversion
			\item $G_{\exp}$ itself has $\sim 0.02\%$ measurement uncertainty
			\item Dimensional factors: $[E^{-1}] \to [E^{-2}] \to [\text{m}^3\text{kg}^{-1}\text{s}^{-2}]$
		\end{itemize}
	\end{tcolorbox}
	
	\section{Why No Fractal Corrections?}
	
	\subsection{Ratio Geometry vs. Absolute Scales}
	
	\begin{theorem}
		\textbf{Ratio Invariance of the Koide Formula}
		
		The Koide formula works exclusively with mass ratios:
		\begin{equation}
			Q = \frac{m_e + m_\mu + m_\tau}{(\sqrt{m_e} + \sqrt{m_\mu} + \sqrt{m_\tau})^2}
		\end{equation}
		
		Since all masses $m_i = r_i \xi^{p_i} v$, the $\xi$-factors partially cancel:
		\begin{equation}
			Q \propto \frac{\xi^{p_1} + \xi^{p_2} + \xi^{p_3}}{(\xi^{p_1/2} + \xi^{p_2/2} + \xi^{p_3/2})^2}
		\end{equation}
		
		The result depends only on the exponent differences:
		\begin{equation}
			\Delta p_{12} = p_1 - p_2, \quad \Delta p_{23} = p_2 - p_3
		\end{equation}
	\end{theorem}
	
	\subsection{Fractal Corrections Only for Absolute Scales}
	
	\begin{table}[h]
		\centering
		\adjustbox{max width=\textwidth}{%
\begin{tabular}{lcc}
			\toprule
			\textbf{Constant} & \textbf{Type} & \textbf{Fractal Correction?} \\
			\midrule
			$Q$ (Koide) & Ratio & \textbf{NO} \\
			$m_p/m_e$ & Ratio & \textbf{NO} \\
			$\alpha$ & Absolute with Scale & \textbf{MINIMAL} \\
			$G$ & Absolute with SI & \textbf{YES} \\
			\bottomrule
		\end{tabular}%
}
		\caption{Necessity of Fractal Corrections}
	\end{table}
	
	% NEW SECTION: Extensions of the Koide Formula

	\section{Unified Theory of Fundamental Constants}
	
	\begin{folgerung}
		\textbf{All three fundamental constants arise from $\xi$:}
		
		\begin{align}
			\text{Koide: } & Q = f_1(\xi^{p_i - p_j}) = \frac{2}{3} \quad &&\text{(Error: } 0.16\%) \\
			\text{Fine Structure: } & \alpha = \xi \cdot E_0^2 = \frac{1}{137.036} \quad &&\text{(Error: } 0.0005\%) \\
			\text{Gravitation: } & G = f_2(\xi, \ell_P) = 6.674 \times 10^{-11} \quad &&\text{(Error: } 0.003\%)
		\end{align}
		
		\textit{(Corrected on 29~Sept.~2026: previously error $0.5\%$ -- with the corpus formula the deviation is $+0.003\%$.)}
		\textit{(Corrected on 30~Sept.~2026: Koide row previously error $0.00003\%$ -- with $m_i = r_i \xi^{p_i} v$, $Q = 0.667742$, i.e.\ $+0.161\%$ above $2/3$.)}
		
		The different accuracies reflect the complexity of the $\xi$-manifestation.
	\end{folgerung}
	
	\subsection{Fundamental Relationship}
	
	The T0 theory reveals a deep connection:
	
	\begin{equation}
		\boxed{\xi \xrightarrow{\text{Ratios}} Q = \frac{2}{3} \xrightarrow{\text{Scale}} \alpha \xrightarrow{\text{SI Units}} G}
	\end{equation}
	
	Each level adds a layer of complexity:
	\begin{itemize}
		\item \textbf{Koide:} Pure Geometry
		\item \textbf{$\alpha$:} Geometry + Energy Scale
		\item \textbf{$G$:} Geometry + Energy Scale + Space-Time Metric
	\end{itemize}
	
% ============================================================
% Addendum 7 September 2026 — P1: Two precision values
% Source: Doc. 190-Archive-1 (P1)
% Source document remains unchanged (append-only, cf. R50).
% ============================================================

\section*{Addendum P1: Clarification of Precision Values (7~September 2026)}

Two different precision values for the Koide relation appear in this document:
$\Delta Q < 0.00003\,\%$ (internal calculation) and $0.001\,\%$ (external presentation).

\textbf{Both are correct} (P1, Doc.~190-Archive-1) --- they describe different
things: $\Delta Q < 0.00003\,\%$ is the direct numerical result of the
calculation using PDG masses; $0.001\,\%$ is the rounded figure for external
publication, which incorporates the measurement uncertainty of the input masses.
The respective context should be observed when reading.

\textit{(Note of 30~Sept.~2026: The value $0.00003\,\%$ does not follow from any PDG masses. With $m_e = 0.51099895$~MeV, $m_\mu = 105.6583755$~MeV and $m_\tau = 1776.93 \pm 0.09$~MeV, $Q - 2/3 = -2.2\times10^{-6}$, relative $-0.00033\,\%$; the $m_\tau$ uncertainty gives $\pm 0.0008\,\%$. The order $0.001\,\%$ thus corresponds to the measurement accuracy; $0.00003\,\%$ as a direct computational result is not tenable. The deviation computed with the T0 masses $m_i = r_i\xi^{p_i}v$ is $+0.161\,\%$.)}
